%%%%%%%%%%%%%%%%%%%%%%%%%%%%%%%%%%%%%%%%%%%%%%%%%%%%%%%%%%%%%%%%%%%%%%%%
% Plantilla TFG/TFM
% Escuela Politécnica Superior de la Universidad de Alicante
% Realizado por: Jose Manuel Requena Plens
% Contacto: info@jmrplens.com / Telegram:@jmrplens
%%%%%%%%%%%%%%%%%%%%%%%%%%%%%%%%%%%%%%%%%%%%%%%%%%%%%%%%%%%%%%%%%%%%%%%%

\chapter{Conclusiones}
\label{conclusiones}

El desarrollo de este Trabajo Final de Máster ha permitido abordar de forma integral y la implementación de un sistema de monitorización de red accesible vía web, aplicado a un entorno de dispositivos de red reales o simulados. Más allá de los resultados técnicos obtenidos, el proyecto ha supuesto un proceso de aprendizaje progresivo en el que se han afrontado y resuelto diversas dificultades propias de un sistema distribuido y dependiente de la conectividad entre múltiples componentes.

Uno de los aspectos más relevantes del desarrollo ha sido la gestión de los problemas de conexión entre los distintos elementos del sistema, especialmente durante las fases iniciales del proyecto. La comunicación entre el backend y los routers, así como la interacción con el entorno de simulación, planteó diversas limitaciones técnicas que obligaron a replantear el entorno de trabajo y a adoptar soluciones más estables. Estas dificultades, lejos de constituir un obstáculo insalvable, permitieron comprender en mayor profundidad la importancia de la configuración del entrono, la fiabilidad de alas conexiones y la necesidad de validar cada componente de forma incremental antes de avanzar en el desarrollo.

Superadas estas limitaciones iniciales, el proyecto evolucionó hacia una arquitectura más robusta, en la que el backend actúa como núcle central de comunicación y procesamiento, separando la lógica de negocio de la capa de presentación. Este desglose ha resultado clave para facilitar tanto el desarrollo como la posterior validación del sistema, permitiendo integrar de forma ordenada las distintas funcionalidades de monitorización y gestión de red.

El sistema desarrollado permite obtener y visualizar información relevante de routers Cisco, incluyendo el estado de las interfaces, el tráfico, los contadores de errores y el uso de recursos. La incorporación de unsistema de detección de errores y notificación automática mediante informes enviados por correo electrónico aporta un enfoque más proactivo a la supervisión de la red, alejándose de una monitorización puramente pasiva basada únicamente en la consulta manual de datos.

Desde el punto de vista del usuario, la aplicación ofrece una interfaz web clara y estructurada que facilita el acceso a la información de monitorización y a las acciones básicas de gestión. Las pruebas realizadas han confirmado que el sistema responde correctamente a los escenarios de uso previstos, validando la correcta integración entre frontend, backend y dispositivos de red, así como la estabilidad del conjunto una vez consolidado el entorno de trabajo.

A nivel académico y personal, la realización de este trabajo ha permitido consolidad conocimientos fundamentales en el ámbito de las redes de comunicaciones, al tiempo que se han reforzado competencias relacionadas con el diseño de arquitecturas software, la resolución de problemas técnicos reales y la toma de decisiones durante el desarrollo de sistemas complejos. En este sentido, el proyecto no solo ha dado lugar a un sistema funcional, sino que ha contribuido de manera significativa al proceso de apredizaje y maduración técnica asociado a la finalización del máster.



